\section{Fitting a check at the answer}\label{sec:setting}

An example is an input $x$ on a set $V$ of $n$ items, together with its answer $y$. The input holds numbers on items
(keys, times, weights), relations on pairs of items (an edge, two letters that are the same) and numbers on pairs (edge
weights); numbers enter only through comparisons. The answer is a set of bits on items or on pairs. A pointer answer,
such as the parent tree of a search, is the pair bit $\pi(u,v)$, ``$u$ is the parent of $v$'', with $\pi(v,v)$ for a
root. Our state is the one of our cold layer\cite{mohsenvand2026zero}: bits and small whole numbers for every item and
every ordered pair, the answer included.

\subsection{Completing the answer}

An answer often hides what decided it. A parent tree does not show how deep each node lies, and the depth is what tells
a breadth-first parent from any other neighbour. We therefore add registers by two generic operations, closure and
counting. For a pointer answer,
\begin{equation}
\mathrm{anc}_y = \pi^{+}, \qquad \mathrm{depth}_y(v) = \#\{u \neq v : \mathrm{anc}_y(u,v)\},
\label{eq:completion}
\end{equation}
where $\pi^{+}$ is the transitive closure of the pointer, counted from the source: a node that the pointer does not
connect to the source receives the value $\notyet$, read ``not yet''. Comparisons treat $\notyet$ as later than every
number: for a number $a$, $a < \notyet$ and $\notyet > a$ hold (and so do $a \leq \notyet$ and $\notyet \geq a$), and
every other comparison with $\notyet$, including $\notyet$ against $\notyet$, is false. The other completions are of the
same two kinds. Three read the answer: the depth, the sum of the weights along the pointer, and, for the longest common
subsequence, the number of match arrows along the answer's own arrows. Two read the input: the bottleneck closure of the
edge weights (for each pair, the least possible largest weight on a path between them) and, for the string matchers, the
closure of ``the text and the pattern agree from here to the pattern's end''. We write $y^{+} = (y, \kappa(y))$ for the
completed answer. No step of any algorithm enters $\kappa$, and the same few operations serve every task.

\subsection{The check}

For a target bit on a pair $(u,v)$, the check is a conjunction of conditions drawn from one fixed family $\Phi$, built
from the input's types alone:
\begin{itemize}
\setlength{\itemsep}{0pt}
\item the pair's own comparisons: $r(u)$ against $r(v)$ ($<$, $\leq$, $=$, $\geq$, $>$) for item registers $r$ of the input
or the completion; one pair number against another on the same pair ($<$, $=$, $>$), such as an edge weight against the
bottleneck; the input relations between $u$ and $v$; and the item bits of $u$ and $v$;
\item witness tests through a third item, $\exists k \notin \{u,v\} : A(u,k) \wedge B(k,v)$, where a leg test is an
own-style test on the pair it joins, optionally a relation or pair-number comparison conjoined with a register
comparison, optionally with a bit of $k$; and pair-number tests through a third item, $\exists k \notin \{u,v\} :
B(k,v) \wedge N_1(k,v) \mathbin{\mathrm{op}} N_2(u,v)$;
\item for an item target $i$, comparisons of two counts, $\#\{k : A(k,i)\} < \#\{k : B(k,i)\}$ or $=$.
\end{itemize}
Each condition can enter a rule negated, and since $\notyet$ fails every comparison it enters, a negated comparison and
its strict counterpart agree on numbers but differ on $\notyet$. With $C \subset \Phi$ the conditions the fit chooses,
\begin{equation}
R(u,v) = \bigwedge_{\varphi \in C} \varphi(x, y^{+})(u,v), \qquad
\pi(v) = \operatorname*{arg\,max}_{u \,:\, R(u,v)} s(u,v),
\label{eq:check}
\end{equation}
where a score $s$ breaks ties among the candidates that pass; the self-pointer $u = v$, ``$v$ has no parent'', competes
with the others by its own rule and score. These are the reads of the cold layer (Fig.~\ref{fig:arch}): a witness test is
its existence read, a count comparison one of its verdict bits\cite{mohsenvand2026zero}. Nothing in $\Phi$ names a
position or a size, so a rule fitted on graphs of 16 nodes applies, unchanged, to graphs of a thousand.

\subsection{Fitting at the answer}

As CRAM does, we predict each bit of a training answer from the rest of the completed answer and the input, with no
loop. Because the answers are exact, we fit the rule of each bit as a conjunction, in two steps. We first keep every own
condition that held on all training positives and failed on at least $\max(10, 0.5\%)$ of the training rows,
\begin{equation}
K = \{\varphi \in \Phi_{\mathrm{own}} : \varphi(u,v) = 1 \text{ for every positive } (u,v) \text{ of every training answer}\},
\label{eq:kept}
\end{equation}
where $\Phi_{\mathrm{own}}$ holds the pair's own conditions. We then add conditions greedily, each true on every positive:
at each step, the one that rejects the most training negatives that still pass, until none passes. This is the classical
learner of a conjunction\cite{haussler1988quantifying}, and $C = K \cup G$, with $G$ the conditions it adds. A logistic
regression with an $\ell_1$ penalty over all of $\Phi$, the pseudolikelihood of each bit given the rest, supplies the
score $s$. A target's own completed value enters its register comparisons only as ``the other item is below me'' ($r(u) < r(v)$
with $r(v)$ the target's), so that it cannot read off the answer it is meant to produce. For breadth-first search, for example, the fit
chooses among \nm{bfs.clrs.features} candidate conditions from \nm{clrs.train.inputs} training answers.

The kept set $K$ is a safeguard. A condition true on every training positive cannot help to reject a negative that the
other conditions already reject, so a sparse or greedy fit has no reason to include it. One such condition is that the
parent lies below the child, $\mathrm{depth}(u) < \mathrm{depth}(v)$. It is true at every finished answer, and it is
the condition that matters on the way to one (Section~\ref{sec:theory}). It also makes ``$u$ is reached'' redundant,
since $\mathrm{depth}(u) < \mathrm{depth}(v)$ is false whenever $\mathrm{depth}(u) = \notyet$.

On puzzles, where every cell holds one of several values, we follow CRAM more closely (Methods). Pseudolikelihood with
an $\ell_1$ penalty chooses which relations between cells carry a cost, from a vocabulary of coordinate relations such
as ``same row'' and ``same block''. For each relation it keeps, we record, over all training answers, the range of the
number of related cells that hold a value. The check passes a board when every such count lies in its range.

\subsection{Running the check cold}

The cold run starts from the empty answer, with no pointer, no selection and every register at $\notyet$, and applies
the check to every item at once, until nothing changes:
\begin{equation}
s_0 = \varnothing, \qquad s_{t+1} = R\bigl(x,\, s_t,\, \kappa(s_t)\bigr), \qquad \text{until } s_{t+1} = s_t .
\label{eq:run}
\end{equation}
Each loop recomputes the completion from the current partial answer by the same operations as at the answer, and writes
a pointer only where a candidate passes; a pointer may change from loop to loop. A node no candidate reaches stays
without a pointer, and at rest it points to itself, the library's ``none''. On puzzles the cold run keeps, for each cell,
the values not yet ruled out, and removes a value when some kept count range can no longer be met. This is the bounds
propagation of constraint programming\cite{rossi2006handbook}, applied to fitted ranges. Where cells stay open, a
depth-first search tries the open cell with the fewest values left, runs the cold stage after each choice and undoes it
on a contradiction. The fitted check is the stopping test in both stages.
